% Options for packages loaded elsewhere
\PassOptionsToPackage{unicode}{hyperref}
\PassOptionsToPackage{hyphens}{url}
\PassOptionsToPackage{dvipsnames,svgnames,x11names}{xcolor}
\documentclass[
  10pt,
  fontset=fandol]{ctexart}
\usepackage{xcolor}
\usepackage[margin=16mm]{geometry}
\usepackage{amsmath,amssymb}
\setcounter{secnumdepth}{-\maxdimen} % remove section numbering
\usepackage{iftex}
\ifPDFTeX
  \usepackage[T1]{fontenc}
  \usepackage[utf8]{inputenc}
  \usepackage{textcomp} % provide euro and other symbols
\else % if luatex or xetex
  \usepackage{unicode-math} % this also loads fontspec
  \defaultfontfeatures{Scale=MatchLowercase}
  \defaultfontfeatures[\rmfamily]{Ligatures=TeX,Scale=1}
\fi
\usepackage{lmodern}
\ifPDFTeX\else
  % xetex/luatex font selection
\fi
% Use upquote if available, for straight quotes in verbatim environments
\IfFileExists{upquote.sty}{\usepackage{upquote}}{}
\IfFileExists{microtype.sty}{% use microtype if available
  \usepackage[]{microtype}
  \UseMicrotypeSet[protrusion]{basicmath} % disable protrusion for tt fonts
}{}
\makeatletter
\@ifundefined{KOMAClassName}{% if non-KOMA class
  \IfFileExists{parskip.sty}{%
    \usepackage{parskip}
  }{% else
    \setlength{\parindent}{0pt}
    \setlength{\parskip}{6pt plus 2pt minus 1pt}}
}{% if KOMA class
  \KOMAoptions{parskip=half}}
\makeatother
\usepackage{longtable,booktabs,array}
\newcounter{none} % for unnumbered tables
\usepackage{calc} % for calculating minipage widths
% Correct order of tables after \paragraph or \subparagraph
\usepackage{etoolbox}
\makeatletter
\patchcmd\longtable{\par}{\if@noskipsec\mbox{}\fi\par}{}{}
\makeatother
% Allow footnotes in longtable head/foot
\IfFileExists{footnotehyper.sty}{\usepackage{footnotehyper}}{\usepackage{footnote}}
\makesavenoteenv{longtable}
\setlength{\emergencystretch}{3em} % prevent overfull lines
\providecommand{\tightlist}{%
  \setlength{\itemsep}{0pt}\setlength{\parskip}{0pt}}
\usepackage{amsmath,amssymb,mathtools,needspace}
\xeCJKDeclareCharClass{CJK}{"2032,"0400->"04FF,"2160->"216B,"2460->"2473,"25A0,"25C7,"25CB,"25CF}
\IfFontExistsTF{Arial Unicode MS}{\xeCJKsetup{AutoFallBack=true}\setCJKfallbackfamilyfont{\CJKrmdefault}{Arial Unicode MS}}{}
\setcounter{secnumdepth}{0}
\setlength{\emergencystretch}{2em}
\allowdisplaybreaks
\usepackage{bookmark}
\IfFileExists{xurl.sty}{\usepackage{xurl}}{} % add URL line breaks if available
\urlstyle{same}
\hypersetup{
  pdftitle={解线性方程组的超松弛迭代法},
  colorlinks=true,
  linkcolor={Maroon},
  filecolor={Maroon},
  citecolor={Blue},
  urlcolor={Blue},
  pdfcreator={LaTeX via pandoc}}

\title{解线性方程组的超松弛迭代法}
\author{}
\date{}

\begin{document}
\maketitle

\subsection{8.4　解线性方程组的超松弛迭代法}\label{ux89e3ux7ebfux6027ux65b9ux7a0bux7ec4ux7684ux8d85ux677eux5f1bux8fedux4ee3ux6cd5}

逐次超松弛迭代法（successive over relaxation method，简称 SOR 方法）是 Gauss-Seidel 方法的一种加速方法，是解大型稀疏矩阵方程组的有效方法之一，它具有计算公式简单，程序设计容易，占用计算机内存较少等优点，但需要选择好的加速因子（即最佳松弛因子）。

设有方程组

\[
Ax=b,
\tag{8.4.1}
\]

其中 \(A\in\mathbb R^{n\times n}\) 为非奇异矩阵，且设 \(a_{ii}\neq0\)（\(i=1,2,\cdots,n\)），分解 \(A\) 为

\[
A=D-L-U.
\tag{8.4.2}
\]

设已知第 \(k\) 次迭代向量 \(x^{(k)}\)，及第 \(k+1\) 次迭代向量 \(x^{(k+1)}\) 的分量 \(x_j^{(k+1)}\)（\(j=1,2,\cdots,i-1\)），要求计算分量 \(x_i^{(k+1)}\)。

首先用 Gauss-Seidel 迭代法定义辅助量

\[
\tilde{x}_i^{(k+1)}=\frac1{a_{ii}}\left(b_i-\sum_{j=1}^{i-1}a_{ij}x_j^{(k+1)}-\sum_{j=i+1}^{n}a_{ij}x_j^{(k)}\right)\quad(i=1,2,\cdots,n),
\tag{8.4.3}
\]

再把 \(x_i^{(k+1)}\) 取为 \(x_i^{(k)}\) 与 \(\tilde{x}_i^{(k+1)}\) 某个平均值（即加权平均），即

\[
x_i^{(k+1)}=(1-\omega)x_i^{(k)}+\omega\tilde{x}_i^{(k+1)}
=x_i^{(k)}+\omega(\tilde{x}_i^{(k+1)}-x_i^{(k)}).
\tag{8.4.4}
\]

用式（8.4.3）代入式（8.4.4）即得到解方程组 \(Ax=b\) 的逐次超松弛迭代公式

\[
\begin{cases}
x_i^{(k+1)}=x_i^{(k)}+\dfrac{\omega}{a_{ii}}\left(b_i-\displaystyle\sum_{j=1}^{i-1}a_{ij}x_j^{(k+1)}-\displaystyle\sum_{j=i}^n a_{ij}x_j^{(k)}\right),\\[5pt]
x^{(k)}=(x_1^{(k)},x_2^{(k)},\cdots,x_n^{(k)})^{\mathrm T}\quad(k=0,1,\cdots;\ i=1,2,\cdots,n).
\end{cases}
\tag{8.4.5}
\]

其中 \(\omega\) 称为\textbf{松弛因子}，或写为

\[
\begin{cases}
x_i^{(k+1)}=x_i^{(k)}+\Delta x_i\quad(k=0,1,\cdots;\ i=1,2,\cdots,n),\\[5pt]
\Delta x_i=\dfrac{\omega}{a_{ii}}\left(b_i-\displaystyle\sum_{j=1}^{i-1}a_{ij}x_j^{(k+1)}-\displaystyle\sum_{j=i}^n a_{ij}x_j^{(k)}\right).
\end{cases}
\tag{8.4.6}
\]

显然，当 \(\omega=1\) 时，解式（8.4.1）的 SOR 方法就是 Gauss-Seidel 迭代法。

在 SOR 方法中，迭代一次主要的运算量是计算一次矩阵与向量的乘法。由式（8.4.5）可知，在计算机上应用 SOR 方法解方程组时只需一组工作单元，以便存放近似

\href{../../source-images/pdf-227.jpeg}{原页图像}

解。在用计算机计算时，可用 \(|p_0|=\max\limits_i|\Delta x_i|=\max\limits_{1\leqslant i\leqslant n}|x_i^{(k+1)}-x_i^{(k)}|<\varepsilon\) 控制迭代终止。

当 \(\omega<1\) 时，称式（8.4.5）为\textbf{低松弛法}；当 \(\omega>1\) 时，称式（8.4.5）为\textbf{超松弛法}。

\textbf{例 8.11}　用 SOR 方法解方程组

\[
\begin{bmatrix}-4&1&1&1\\1&-4&1&1\\1&1&-4&1\\1&1&1&-4\end{bmatrix}
\begin{bmatrix}x_1\\x_2\\x_3\\x_4\end{bmatrix}
=\begin{bmatrix}1\\1\\1\\1\end{bmatrix},
\]

它的精确解为 \(x^*=(-1,-1,-1,-1)^{\mathrm T}\)。

\textbf{解}　取 \(x^{(0)}=0\)，迭代公式为

\[
\begin{cases}
x_1^{(k+1)}=x_1^{(k)}-\omega(1+4x_1^{(k)}-x_2^{(k)}-x_3^{(k)}-x_4^{(k)})/4,\\
x_2^{(k+1)}=x_2^{(k)}-\omega(1-x_1^{(k+1)}+4x_2^{(k)}-x_3^{(k)}-x_4^{(k)})/4,\\
x_3^{(k+1)}=x_3^{(k)}-\omega(1-x_1^{(k+1)}-x_2^{(k+1)}+4x_3^{(k)}-x_4^{(k)})/4,\\
x_4^{(k+1)}=x_4^{(k)}-\omega(1-x_1^{(k+1)}-x_2^{(k+1)}-x_3^{(k+1)}+4x_4^{(k)})/4.
\end{cases}
\]

取 \(\omega=1.3\)，第 11 次迭代结果为

\[
x^{(11)}=(-0.999\,996\,46,\ -1.000\,003\,10,\ -0.999\,999\,53,\ -0.999\,999\,12)^{\mathrm T},
\]

\[
\|\varepsilon^{(11)}\|_2\leqslant0.46\times10^{-5}.
\]

对 \(\omega\) 取其他值，迭代次数如表 8.1 所示。从此例可以看到，松弛因子选择得好，会使 SOR 方法的收敛大大加速。本例中，\(\omega=1.3\) 是最佳松弛因子。

\Needspace{19\baselineskip}

\textbf{表 8.1}

{\def\LTcaptype{none} % do not increment counter
\begin{longtable}[]{@{}lr@{}}
\toprule\noalign{}
松弛因子 \(\omega\) & 满足误差 \(\|x^{(k)}-x^*\|_2<10^{-5}\) 的迭代次数 \\
\midrule\noalign{}
\endhead
\bottomrule\noalign{}
\endlastfoot
1.0 & 22 \\
1.1 & 17 \\
1.2 & 12 \\
\(\underline{1.3}\) & 11（最少迭代次数） \\
1.4 & 14 \\
1.5 & 17 \\
1.6 & 23 \\
1.7 & 33 \\
1.8 & 53 \\
1.9 & 109 \\
\end{longtable}
}

下面写出 SOR 迭代公式的矩阵形式。迭代公式（8.4.5）亦可写为

\[
a_{ii}x_i^{(k+1)}=(1-\omega)a_{ii}x_i^{(k)}+\omega\left(b_i-\sum_{j=1}^{i-1}a_{ij}x_j^{(k+1)}-\sum_{j=i+1}^n a_{ij}x_j^{(k)}\right)\quad(i=1,2,\cdots,n).
\tag{8.4.7}
\]

用式 \(A=D-L-U\)，则

\[
Dx^{(k+1)}=\omega(b+Lx^{(k+1)}+Ux^{(k)})+(1-\omega)Dx^{(k)},
\]

即

\[
(D-\omega L)x^{(k+1)}=((1-\omega)D+\omega U)x^{(k)}+\omega b.
\]

显然对于任何一个 \(\omega\) 值，\(D-\omega L\) 非奇异（由设 \(a_{ii}\neq0;\ i=1,2,\cdots,n\)），于是

\[
x^{(k+1)}=(D-\omega L)^{-1}[(1-\omega)D+\omega U]x^{(k)}+\omega(D-\omega L)^{-1}b.
\]

这就是说，解式（8.2.1）（设 \(a_{ii}\neq0;\ i=1,2,\cdots,n\)）的 SOR 方法迭代公式为

\[
x^{(k+1)}=L_\omega x^{(k)}+f.
\tag{8.4.8}
\]

\begin{quote}
校注（agent 补充）：原页第 11 次迭代向量及误差界均按原样保留。按所印四个分量计算，\(\|x^{(11)}-x^*\|_2\approx4.8100832\times10^{-6}\)，超过所印上界 \(0.46\times10^{-5}\)。以 \(\omega=13/10\) 从零向量作 11 次有理数迭代，所得向量约为 \((-0.9999966722,-1.0000028727,-0.9999995354,-0.9999991925)^{\mathrm T}\)，误差约为 \(4.4938646\times10^{-6}\)；因此本页所印向量与误差界存在数值不一致，并非 OCR 错误。
\end{quote}

\href{../../source-images/pdf-228.jpeg}{原页图像}

其中

\[
L_\omega=(D-\omega L)^{-1}[(1-\omega)D+\omega U],\qquad f=\omega(D-\omega L)^{-1}b.
\]

矩阵 \(L_\omega\) 称为 SOR 方法的迭代矩阵。这说明对式（8.2.1）应用 SOR 方法相当于对方程组 \(x=L_\omega x+f\) 应用一般迭代法。于是关于一般迭代法的理论可应用到式（8.4.8），得到下述定理。

\textbf{定理 8.8}　设有线性方程组（8.4.1），且 \(a_{ii}\neq0\)（\(i=1,2,\cdots,n\)），则解方程组的 SOR 方法收敛的充要条件是

\[
\rho(L_\omega)<1.
\]

引进超松弛迭代法的想法是希望能选择松弛因子 \(\omega\) 使得迭代过程式（8.4.5）收敛较快，也就是应选择因子 \(\omega\) 使 \(\rho(L_\omega)=\min\limits_\omega\)。

下面研究，对于一般方程组（8.2.1）（\(a_{ii}\neq0;\ i=1,2,\cdots,n\)），松弛因子 \(\omega\) 在什么范围内取值，SOR 方法才可能收敛。现给出 SOR 方法收敛的必要条件。

\textbf{定理 8.9}　设解式（8.2.1）（\(a_{ii}\neq0;\ i=1,2,\cdots,n\)）的 SOR 方法收敛，则

\[
0<\omega<2.
\]

\textbf{证明}　由设 SOR 方法收敛，根据定理 8.8，\(\rho(L_\omega)<1\)。

设 \(L_\omega\) 的特征值为 \(\lambda_1,\lambda_2,\cdots,\lambda_n\)，则

\[
|\det L_\omega|=|\lambda_1\lambda_2\cdots\lambda_n|\leqslant(\rho(L_\omega))^n,\qquad
|\det L_\omega|^{1/n}\leqslant\rho(L_\omega)<1.
\]

而

\[
\det L_\omega=\det((D-\omega L)^{-1})\cdot\det((1-\omega)D+\omega U)=(1-\omega)^n,
\]

所以

\[
|1-\omega|<1.
\]

该定理说明对于解一般方程组（8.2.1）（\(a_{ii}\neq0;\ i=1,2,\cdots,n\)），SOR 方法只有取松弛因子 \(\omega\) 在 \((0,2)\) 范围内才能收敛。而当 \(A\) 是对称正定矩阵时，若 \(\omega\) 满足 \(0<\omega<2\)，则 SOR 方法一定收敛。

\textbf{定理 8.10}　如果 \(A\) 为对称正定矩阵，且 \(0<\omega<2\)，则解式（8.4.1）的 SOR 方法收敛。

\textbf{证明}　在上述假定下，若能证明 \(|\lambda|<1\)，那么定理得证（其中 \(\lambda\) 为 \(L_\omega\) 的任一特征值）。

事实上，设 \(y\) 为对应 \(\lambda\) 的 \(L_\omega\) 的特征向量，即

\[
L_\omega y=\lambda y,\qquad y=(y_1,y_2,\cdots,y_n)^{\mathrm T}\neq0,
\]

\[
(D-\omega L)^{-1}[(1-\omega)D+\omega U]y=\lambda y,
\]

亦即

\[
[(1-\omega)D+\omega U]y=\lambda(D-\omega L)y.
\]

为了找出 \(\lambda\) 的表达式，考虑数量积

\[
(((1-\omega)D+\omega U)y,y)=\lambda((D-\omega L)y,y),
\]

则

\[
\lambda=\frac{(Dy,y)-\omega(Dy,y)+\omega(Uy,y)}{(Dy,y)-\omega(Ly,y)},
\]

\href{../../source-images/pdf-229.jpeg}{原页图像}

显然

\[
(Dy,y)=\sum_{i=1}^n a_{ii}|y_i|^2\equiv\sigma>0,
\tag{8.4.9}
\]

记 \(-(Ly,y)=\alpha+\mathrm i\beta\)，由于 \(A=A^{\mathrm T}\)，所以 \(U=L^{\mathrm T}\)。

\[
-(Uy,y)=-(y,Ly)=-\overline{(Ly,y)}=\alpha-\mathrm i\beta,
\]

\[
0<(Ay,y)=((D-L-U)y,y)=\sigma+2\alpha,
\tag{8.4.10}
\]

所以

\[
\lambda=\frac{(\sigma-\omega\sigma-\alpha\omega)+\mathrm i\omega\beta}{(\sigma+\alpha\omega)+\mathrm i\omega\beta},
\]

从而

\[
|\lambda|^2=\frac{(\sigma-\omega\sigma-\alpha\omega)^2+\omega^2\beta^2}{(\sigma+\alpha\omega)^2+\omega^2\beta^2}.
\]

当 \(0<\omega<2\) 时，利用式（8.4.9）、式（8.4.10），有

\[
(\sigma-\omega\sigma-\alpha\omega)^2-(\sigma+\alpha\omega)^2
=\omega\sigma(\sigma+2\sigma)(\omega-2)<0,
\]

即 \(L_\omega\) 的任一特征值满足 \(|\lambda|<1\)，故 SOR 方法收敛。注意，当 \(0<\omega<2\) 时，可以证明

\[
(\sigma+2\omega)^2+\omega^2\beta^2\neq0.
\]

最佳松弛因子理论是由 Young（1950 年）针对一类椭圆型微分方程数值解得到的代数方程组 \(Ax=b\)（具有所谓性质 A 和相容次序）所建立的理论。他给出了最佳松弛因子公式

\[
\omega_{\mathrm{opt}}=\frac2{1+\sqrt{1-\rho^2(B_0)}},
\]

其中 \(\rho(B_0)\) 是 Jacobi 方法迭代矩阵 \(B_0\) 的谱半径。

一般来说，在实际应用中计算 \(\rho(B_0)\) 较困难，对某些微分方程数值解问题可考虑用第 9 章求近似值的方法，亦可由计算实践摸索出（近似）最佳松弛因子。

\textbf{算法}　本算法用 SOR 方法解式（8.2.1），其中 \(A\) 为对称正定矩阵。数组 \(x\) 为一组工作单元，开始存放初始向量，然后存放近似值解 \(x^{(k)}\)，最后存放结果。用

\[
|p_0|=\max_{1\leqslant i\leqslant n}|\Delta x_i|=\max_{1\leqslant i\leqslant n}|x_i^{(k+1)}-x_i^{(k)}|<\varepsilon\quad\text{（精度要求）}
\]

控制迭代终止。\(k\) 表示迭代次数，可以不用。

\textbf{步 1}　\(k\leftarrow0\)。

\textbf{步 2}　\(x_i\leftarrow0\)（\(i=1,2,\cdots,n\)）。

\textbf{步 3}　\(k\leftarrow k+1\)。

\textbf{步 4}　\(p_0\leftarrow0\)。

\textbf{步 5}　对于 \(i=1,2,\cdots,n\)，有

（1）\(p\leftarrow\Delta x_i=\omega\left(b_i-\displaystyle\sum_{j=1}^{i-1}a_{ij}x_j-\displaystyle\sum_{j=i+1}^n a_{ij}x_j\right)/a_{ii}\)；

（2）如果 \(|p|>|p_0|\)，则 \(p_0\leftarrow p\)；

（3）\(x_i\leftarrow x_i+p\)。

\textbf{步 6}　输出 \(p_0\)。

\begin{quote}
校注（agent 补充）：本页存在以下三处疑似原书排印错误，均已放大核对并保留原式，并非 OCR 错误。第一，平方差的因子原印 \(\sigma+2\sigma\)，代数展开应得 \(\sigma+2\alpha\)，与式（8.4.10）一致。第二，分母非零说明原印 \((\sigma+2\omega)^2+\omega^2\beta^2\neq0\)，由上文 \(|\lambda|^2\) 表达式，相应分母应为 \((\sigma+\alpha\omega)^2+\omega^2\beta^2\)。第三，算法步 5（1）的第二个求和下限原印 \(j=i+1\)；由于随后执行 \(x_i\leftarrow x_i+p\)，\(p=\Delta x_i\) 应是校正量，对照式（8.4.6），该下限应为 \(j=i\)（或另减去 \(a_{ii}x_i\)），否则漏掉对角项。
\end{quote}

\href{../../source-images/pdf-230.jpeg}{原页图像}

\textbf{步 7}　如果 \(|p_0|>\varepsilon\)，则转步 3。

\textbf{步 8}　输出结果 \(x,k\)。

\subsection{本节覆盖原页的校注记录}\label{ux672cux8282ux8986ux76d6ux539fux9875ux7684ux6821ux6ce8ux8bb0ux5f55}

下列记录按原页关联，含同页相邻小节的校注；计算时同时读取上文校注和完整原页。

\begin{itemize}
\tightlist
\item
  \texttt{ch08-E006}：原书 PDF 页 229
\item
  \texttt{ch08-E007}：原书 PDF 页 229
\item
  \texttt{ch08-E008}：原书 PDF 页 229
\item
  \texttt{ch08-E005}：原书 PDF 页 227
\end{itemize}

\end{document}
